% !TeX root = ../main.tex
% 原始章节：index.Rmd

\chapter*{前言}\label{chap:00-preface}
\addcontentsline{toc}{chapter}{前言}

北京航空航天大学自2008年便开始着重培养学生的计算机系统能力，通过调整与优化人才培养方案，打造了一系列计算机系统领域的重点课程，在计算机组成原理、操作系统、编译原理等核心课程方面取得了突出的教学成果。

2017年以来，学校不断鼓励并组织学生参加了共计八届``龙芯杯''全国大学生计算机系统能力培养大赛，并屡获佳绩。通过这一平台，培养了一大批对计算机系统领域有着浓厚兴趣且具备扎实系统设计与开发能力的优秀学生。其中的大部分学生现已更深入地投身于计算机系统领域研究或相关岗位。

自2019年起，学校进一步组织学生在大赛基础上持续深入开展CPU处理器芯片的自主设计与开发工作，迄今已完成两轮芯片设计与流片工作，主要包括1款MIPS处理器芯片和2款基于LoongArch架构的处理器芯片。两款芯片均一次流片成功，并具有以下几个突出特点：（1）处理器核心经过课程设计、``龙芯杯''竞赛以及本科毕业设计等多个环节的迭代与优化，具备优良的性能与可靠性；（2）具备完善的SoC系统集成并支持多种外设，可进行多种功能的验证和演示，形成较为完整的计算机系统形态；（3）支持启动新版本Linux操作系统，且能稳定运行各类外设驱动；（4）芯片基于自主可控的LoongArch指令集设计并流片，能够为我国高端自主可控处理器芯片人才培养提供宝贵的经验。

经过多年的教学、竞赛和实践活动，我校已基本建立起一套可行设计流程和技术体系。为了传承和推广这些宝贵的经验，我们组织了专门的学生团队，对设计资料和流程进行了系统的整理和总结，以供未来的从业者和相关领域的专业人士参考。

本项工作得到了北京航空航天大学计算机学院的大力支持，并在集成电路学院和软件学院的协同配合下不断丰富与完善。该项目也得到了北京航空航天大学重点教改项目、北京航空航天大学精品劳动教育课程项目、教育部产学合作协同育人项目的直接支持以及龙芯中科公司等单位提供的技术与人力上的帮助。在此，谨向所有支持单位和个人表示诚挚的感谢！
